\documentclass[11pt,a4paper]{article}
\usepackage[utf8]{inputenc}
\usepackage{amsmath,amssymb,amsthm}
\usepackage{physics}
\usepackage{geometry}
\geometry{margin=2.5cm}
\usepackage{hyperref}

\newcommand{\ai}{a_i}
\newcommand{\aj}{a_j}
\newcommand{\Sfour}{S_4}

\title{Monte Carlo move on the smoothed-cube manifold\\
       for the $2{+}4$ $p$-spin model with sparse interactions}
\author{}
\date{}

\begin{document}
\maketitle

%% ============================================================
\section{Model and configuration space}
%% ============================================================

We consider $N$ complex amplitudes $\{a_k\}_{k=1}^{N}$, $a_k\in\mathbb{C}$,
subject to the \emph{smoothed-cube} constraint
\begin{equation}\label{eq:cube}
  \sum_{k=1}^{N} \abs{a_k}^4 = N \,.
\end{equation}
This quartic normalisation defines a smooth, compact sub-manifold of
$\mathbb{C}^N$ that interpolates between the $\ell^4$ unit ball (a
``rounded hypercube'') and the standard sphere $\sum\abs{a_k}^2=\mathrm{const}$.

The Hamiltonian is the sum of a two-body and a four-body interaction,
\begin{equation}\label{eq:H}
  H = H_2 + H_4 \,,
\end{equation}
with
\begin{align}
  H_2 &= -\sum_{i<j} \Re\bigl[ g^{(2)}_{ij}\, a_i\, a_j^* \bigr] \,,
  \label{eq:H2}\\[4pt]
  H_4 &= -\sum_{i<j<k<l}
         \Re\bigl[ g^{(4)}_{ijkl}\, f_{c(ijkl)}(a_i,a_j,a_k,a_l) \bigr] \,,
  \label{eq:H4}
\end{align}
where each ordered quartet $(i<j<k<l)$ is assigned a \emph{unique}
interaction channel $c(ijkl)\in\{0,1,2\}$, and the product $f_c$ depends
on $c$:
\begin{alignat}{2}
  c=0:&\quad f_0 = a_i\, a_j\, a_k^*\, a_l^* \,,
       &\qquad & \text{conjugation on slots } \{3,4\}\,,\notag\\
  c=1:&\quad f_1 = a_i\, a_j^*\, a_k^*\, a_l \,,
       && \text{conjugation on slots } \{2,3\}\,,\notag\\
  c=2:&\quad f_2 = a_i\, a_j^*\, a_k\, a_l^* \,,
       && \text{conjugation on slots } \{2,4\}\,.\notag
\end{alignat}
The coupling $g^{(2)}_{ij}$ is a quenched random variable; $g^{(4)}_{ijkl}$
is likewise random and depends only on the quartet (the spatial overlap
integral is symmetric under permutations of the four indices, so all three
channel candidates share the same coupling value).

The channel $c(ijkl)$ is selected by a \emph{frequency-matching condition}
(FMC).  For each quartet $(i<j<k<l)$ one checks the resonance conditions
\begin{alignat}{2}
  c=0:&\quad \abs{\omega_i+\omega_j-\omega_k-\omega_l} \le \gamma\,,\notag\\
  c=1:&\quad \abs{\omega_i+\omega_l-\omega_j-\omega_k} \le \gamma\,,\notag\\
  c=2:&\quad \abs{\omega_i+\omega_k-\omega_j-\omega_l} \le \gamma\,,\notag
\end{alignat}
in this fixed order; the \emph{first} channel that passes is assigned to
the quartet.  If none passes, the quartet is inactive ($g^{(4)}_{ijkl}$
does not contribute to $H_4$).  In the fully-connected limit
($\gamma\to\infty$, or no FMC) every quartet is assigned channel~$0$.
This ensures that each quartet contributes \emph{at most one} term to the
Hamiltonian, which is a physical requirement.

In the \textbf{sparse} variant the quartic interaction retains only $N$
randomly sub-sampled quartet entries from the $O(N^4)$ FMC-surviving set,
reducing the cost of a single energy evaluation from $O(N^4)$ to $O(N)$.

%% ============================================================
\section{The pair-rotation move}
%% ============================================================

A Monte Carlo (MC) sweep consists of $\lfloor N/2\rfloor$ elementary
\emph{pair-rotation} steps.
Each step selects a uniformly random pair of indices $i\neq j$ and proposes
a new configuration that modifies only $a_i$ and $a_j$ while leaving all
other amplitudes unchanged and preserving the global constraint~\eqref{eq:cube}.

\subsection{Conserved pairwise quantity}

The key observation is that the global constraint $\sum_k\abs{a_k}^4 = N$
is preserved if and only if the \emph{pairwise fourth-moment sum}
\begin{equation}\label{eq:S4}
  \Sfour \;\equiv\; \abs{\ai}^4 + \abs{\aj}^4
\end{equation}
is conserved by the move.  This decouples the pair $(i,j)$ from the rest of
the configuration.

\subsection{Parametrisation and flat measure}\label{sec:param}

We parametrise the pair on the constraint surface $\Sfour = \mathrm{const}$
by three coordinates $(u,\,\varphi_1,\,\varphi_2)$:
\begin{equation}\label{eq:param}
  a_i = \abs{a_i}\, e^{i\varphi_1}\,,\qquad
  a_j = \abs{a_j}\, e^{i\varphi_2}\,,
\end{equation}
with the moduli determined by
\begin{equation}\label{eq:moduli}
  \abs{a_i}^2 = \sqrt{\Sfour}\,\sqrt{1-u}\,,\qquad
  \abs{a_j}^2 = \sqrt{\Sfour}\,\sqrt{u}\,,\qquad
  u\in[0,1]\,.
\end{equation}
One easily verifies that
$\abs{a_i}^4 + \abs{a_j}^4 = \Sfour(1-u) + \Sfour\, u = \Sfour$, so the
constraint is satisfied for all $u$.

The variable $u = \abs{a_j}^4/\Sfour$ is chosen so that the
\emph{Riemannian measure} induced on the constraint surface is flat:
\begin{equation}\label{eq:measure}
  \dd\mu \;\propto\; \dd u\;\dd\varphi_1\;\dd\varphi_2 \,.
\end{equation}
This can be seen via the standard change of variables from the angular
parametrisation $\abs{a_i}^2 = \sqrt{\Sfour}\cos\alpha$,
$\abs{a_j}^2 = \sqrt{\Sfour}\sin\alpha$, where the measure carries a
Jacobian factor $\sin\alpha\cos\alpha\,\dd\alpha$.  Setting
$u = \sin^2\!\alpha$ gives $\dd u = 2\sin\alpha\cos\alpha\,\dd\alpha$, which
precisely cancels the Jacobian and yields~\eqref{eq:measure}.

\medskip\noindent\textbf{Remark.}
An earlier version of the code sampled $\alpha$ uniformly in $[0,\pi/2)$,
which fails to reproduce the measure~\eqref{eq:measure}: the proposal
density $q(\alpha) = 2/\pi$ differs from the constraint-surface measure
$\mu(\alpha)\propto\sin\alpha\cos\alpha$.  For Metropolis--Hastings
correctness one would then need to include the ratio
$\sin\alpha'\cos\alpha'\!/\!\sin\alpha\cos\alpha$ in the acceptance
probability.  The $u$-parametrisation eliminates this complication entirely.

%% ============================================================
\section{Proposal distribution and detailed balance}
%% ============================================================

Given the current spin configuration with values $a_i^{\mathrm{old}}$,
$a_j^{\mathrm{old}}$, the proposal draws
\begin{equation}\label{eq:proposal}
  u \sim \mathrm{Uniform}(0,1)\,,\qquad
  \varphi_1 \sim \mathrm{Uniform}(0,2\pi)\,,\qquad
  \varphi_2 \sim \mathrm{Uniform}(0,2\pi)\,,
\end{equation}
independently.  The proposed amplitudes are then
\begin{equation}\label{eq:new_spins}
  a_i' = \sqrt{\sqrt{\Sfour}\,\sqrt{1-u}}\;
         e^{\,i\varphi_1}\,,\qquad
  a_j' = \sqrt{\sqrt{\Sfour}\,\sqrt{u}}\;
         e^{\,i\varphi_2}\,.
\end{equation}
Note that the proposal~\eqref{eq:proposal} is \emph{independent} of the
current state (it depends only on $\Sfour$, which is a constant of the
move).  The proposal density in the flat coordinates $(u,\varphi_1,\varphi_2)$
is therefore
\[
  q\bigl((u',\varphi_1',\varphi_2') \mid (u,\varphi_1,\varphi_2)\bigr)
  = \frac{1}{1 \cdot 2\pi \cdot 2\pi}
  = \frac{1}{4\pi^2}\,,
\]
which is symmetric:
$q(\mathbf{x}'\mid\mathbf{x}) = q(\mathbf{x}\mid\mathbf{x}')$.

Because both the proposal density and the constraint-surface measure are
uniform in $(u,\varphi_1,\varphi_2)$, the Metropolis--Hastings ratio reduces
to
\begin{equation}\label{eq:acceptance}
  A = \min\!\Bigl(1,\; e^{-\beta\,\Delta E}\Bigr)\,,
\end{equation}
with $\Delta E = H(\text{new}) - H(\text{old})$, and no additional Jacobian
or proposal-ratio correction.

%% ============================================================
\section{Energy difference computation}
%% ============================================================

Because only spins $i$ and $j$ change, the energy difference decomposes
into contributions from $H_2$ and $H_4$.

\subsection{Two-body term}

Defining $\delta_i = a_i' - a_i^{\mathrm{old}}$ and
$\delta_j = a_j' - a_j^{\mathrm{old}}$, the change in $H_2$ can be
written as
\begin{equation}\label{eq:dE2}
  \Delta E_2 = -\Re\Bigl[
    \delta_i \sum_{k\neq i,j} g^{(2)}_{ik}\, a_k^*
  + \delta_j \sum_{k\neq i,j} g^{(2)}_{jk}\, a_k^*
  + g^{(2)}_{ij}\bigl(a_i' a_j'^* - a_i^{\mathrm{old}} a_j^{\mathrm{old}\,*}\bigr)
  \Bigr]\,.
\end{equation}
This requires $O(N)$ work and is computed in parallel over the spin index
$k$ (dense storage of $g^{(2)}$).

\subsection{Four-body term (sparse)}

Recall that each sparse entry $q$ stores a sorted quartet of indices
$(i_1\!<\!i_2\!<\!i_3\!<\!i_4)$, the unique channel label
$c = c(i_1 i_2 i_3 i_4)\in\{0,1,2\}$ selected by FMC, and a
complex coupling $g_q$.  The channel determines a \emph{conjugation mask}
$m_c \in \{0,1\}^4$ specifying which of the four slots carry a complex
conjugation:
\begin{equation}\label{eq:conj}
  m_0 = (0,0,1,1)\,,\qquad
  m_1 = (0,1,1,0)\,,\qquad
  m_2 = (0,1,0,1)\,.
\end{equation}
We define the channel-dependent factor for slot $p$ as
\begin{equation}\label{eq:tilde}
  \tilde a_p^{(c)} \;\equiv\;
    \begin{cases}
      a_{i_p}^{\phantom{*}} & \text{if } (m_c)_p = 0\,,\\[2pt]
      a_{i_p}^{*}           & \text{if } (m_c)_p = 1\,,
    \end{cases}
\end{equation}
so that the contribution of quartet $q$ to $H_4$ reads
\begin{equation}\label{eq:H4q}
  h_q = -\Re\!\Bigl[
    g_q \;\tilde a_1^{(c)}\,\tilde a_2^{(c)}\,
          \tilde a_3^{(c)}\,\tilde a_4^{(c)}
  \Bigr]\,.
\end{equation}

For each sparse entry $q$ we check whether the modified indices $i$ or $j$
(or both) appear among $\{i_1,i_2,i_3,i_4\}$.  If neither appears, the
entry does not contribute to $\Delta E_4$.  Two sub-cases arise.

\paragraph{Case~1: exactly one of $i,j$ appears in the quartet.}
Suppose, without loss of generality, that it is spin $i$ at slot position
$p_0$, and define $\delta_i = a_i' - a_i^{\mathrm{old}}$.  Because
$H_4 \ni -\Re[g_q\,\prod_p \tilde a_p]$ is linear in each slot, the
change due to replacing $a_i\to a_i + \delta_i$ in slot $p_0$ is
\begin{equation}\label{eq:dE4_one}
  \Delta h_q^{(\mathrm{one})} = -\Re\!\Biggl[
    g_q \prod_{p=1}^{4} \tilde b_p^{(c)}
  \Biggr]\,,\qquad
  \tilde b_p^{(c)} =
    \begin{cases}
      \widetilde{\delta_i}^{(c)} & \text{if } p = p_0\,,\\[2pt]
      \tilde a_p^{(c)}           & \text{otherwise}\,,
    \end{cases}
\end{equation}
where $\widetilde{\delta_i}^{(c)}$ carries the same conjugation as
slot $p_0$:
$\widetilde{\delta_i}^{(c)} = \delta_i$ if $(m_c)_{p_0}=0$,
$\widetilde{\delta_i}^{(c)} = \delta_i^*$ if $(m_c)_{p_0}=1$.
Note that this expression is \emph{exact} (not a linearisation), because
$\delta_i$ appears in exactly one slot.

\paragraph{Case~2: both $i$ and $j$ appear in the quartet.}
When both modified spins participate, the cross-term between $\delta_i$ and
$\delta_j$ prevents a simple factorisation.  In this case we recompute the
full product before and after the move:
\begin{equation}\label{eq:dE4_both}
  \Delta h_q^{(\mathrm{both})} = -\Re\!\Biggl[
    g_q\!\Biggl(
      \prod_{p=1}^{4} \tilde a_p^{\prime(c)}
    - \prod_{p=1}^{4} \tilde a_p^{(c)}
    \Biggr)
  \Biggr]\,,
\end{equation}
where $\tilde a_p^{\prime(c)}$ denotes the factor~\eqref{eq:tilde} evaluated
with $a_i \to a_i'$ and $a_j \to a_j'$.

\paragraph{Total $\Delta E_4$.}
Summing over all sparse entries,
\begin{equation}\label{eq:dE4_total}
  \Delta E_4 = \sum_{q=1}^{n_{\mathrm{quartets}}} \Delta h_q \,,\qquad
  \Delta h_q =
    \begin{cases}
      0                              & \text{if } i,j \notin \{i_1,i_2,i_3,i_4\}\,,\\
      \Delta h_q^{(\mathrm{one})}    & \text{if exactly one of } i,j \in \{i_1,i_2,i_3,i_4\}\,,\\
      \Delta h_q^{(\mathrm{both})}   & \text{if both } i,j \in \{i_1,i_2,i_3,i_4\}\,.
    \end{cases}
\end{equation}
The total cost per step is $O(n_{\mathrm{quartets}})$ where
$n_{\mathrm{quartets}} = N$ in the sparse variant, giving an $O(N)$ cost
per MC step and $O(N^2)$ per sweep.

%% ============================================================
\section{Sweep structure and parallelism}
%% ============================================================

One MC sweep performs $\lfloor N/2\rfloor$ sequential pair-rotation steps
(sequential because each step may modify spins that appear in subsequent
proposals).  On GPU, the energy-difference computation for each step is
parallelised across CUDA threads within a single block:
\begin{itemize}
  \item Spins are loaded into shared memory at the start of the sweep.
  \item Thread~0 selects the pair and generates the proposal.
  \item All threads cooperatively compute $\Delta E$ via parallel reduction
    (warp-shuffle followed by a final warp-level sum).
  \item Thread~0 performs the Metropolis test and, if accepted, updates the
    shared-memory copy of the two spins.
\end{itemize}
Each CUDA block handles one replica; in a parallel-tempering run, all
replicas execute their sweeps simultaneously.

%% ============================================================
\section{Summary}
%% ============================================================

The pair-rotation move on the smoothed-cube manifold $\sum\abs{a_k}^4=N$ is
a Metropolis update that:
\begin{enumerate}
  \item Picks a random pair $(i,j)$.
  \item Conserves the pairwise invariant
        $\Sfour = \abs{a_i}^4+\abs{a_j}^4$.
  \item Proposes new moduli and phases by sampling $(u,\varphi_1,\varphi_2)$
        uniformly in $[0,1]\times[0,2\pi)^2$, where $u=\abs{a_j}^4/\Sfour$.
  \item Accepts with probability $\min(1,e^{-\beta\Delta E})$,
        with no Jacobian correction required.
\end{enumerate}
The $u$-parametrisation ensures that the flat measure on the constraint
surface is exactly reproduced by the uniform proposal, guaranteeing detailed
balance without any Jacobian factor in the acceptance ratio.

\end{document}
